\documentclass[11pt]{article}

% ============================================================
%  Overleaf-ready. Compile with pdfLaTeX.
%  Progress report: from the Emerson offline-RL baseline to the
%  PhiBE three-term objective, and why a standard RL optimizer
%  cannot minimize it.
% ============================================================

\usepackage[margin=1in]{geometry}
\usepackage{amsmath,amssymb,amsthm}
\usepackage{bm}
\usepackage{enumitem}
\usepackage{hyperref}
\hypersetup{colorlinks=true,linkcolor=blue,urlcolor=blue,citecolor=blue}

\newcommand{\R}{\mathbb{R}}
\newcommand{\E}{\mathbb{E}}
\newcommand{\grad}{\nabla}
\newcommand{\grads}{\nabla_s}
\newcommand{\hess}{\nabla_s^2}
\newcommand{\tr}{\operatorname{tr}}
\newcommand{\Res}{\mathcal{R}}
\newcommand{\D}{\mathcal{D}}

\title{\textbf{Progress Report: A PhiBE-Style Critic Target for\\
Offline RL in Type-1 Diabetes Control\\[4pt]
\large and Why It Requires a Non-Standard Optimizer}}
\author{Hanzhang Yuan \\ \small Advisor: Prof.\ Yuhua Zhu}
\date{June 24, 2026}

\begin{document}
\maketitle

\begin{abstract}
We start from the offline reinforcement-learning (RL) baseline of
Emerson et al.\ (2023) for blood-glucose control with TD3-BC, and make
\emph{a single modification}: we replace the critic's regression
\emph{target}. Following the continuous-time value function, we Taylor-expand
the expectation term inside the target. This expansion makes the value
function $V$, its state gradient $\grads V$, and its state Hessian $\hess V$
appear \emph{simultaneously} in the objective. The resulting learning
objective therefore has the form $f\big(V_\theta,\ \grads V_\theta,\
\hess V_\theta\big)$, in sharp contrast to a standard RL objective
$f\big(V_\theta\big)$ that depends on $\theta$ only through the network's
output value. We make the objective precise and then analyze why a standard
RL optimizer cannot minimize it. Identifying a suitable optimizer for this
derivative-bearing objective is the focus of the next stage.
\end{abstract}

\section{Setup: the Emerson baseline and the single change we make}

The Emerson et al.\ (2023) baseline trains a TD3-BC critic by regressing
$Q_\theta(s,a)$ onto the standard discrete-time Bellman target
\begin{equation}
\label{eq:std-target}
y^{\mathrm{std}} \;=\; r + \gamma\, Q_{\theta^-}\!\big(s', \pi(s')\big),
\end{equation}
where $\theta^-$ denotes the (frozen) target network. The \emph{only}
change we introduce is to the target $y$. Everything else in TD3-BC---the
actor update, the behavior-cloning term, the replay buffer, and the
target-network mechanism---is left untouched.

\section{The continuous-time value function and its Taylor expansion}

\subsection{Starting point}
We adopt the continuous-time value function
\begin{equation}
\label{eq:ct-value}
V(s_t) \;=\; r\,\Delta t \;+\; e^{-\beta \Delta t}\,
\E\big[\,V(s_{t+1}) \mid s_t\,\big],
\end{equation}
where $\beta>0$ is a continuous-time decay rate (per unit time) and
$\Delta t$ is the sampling interval. (The state-action version,
$Q(s_t,a_t) = r\,\Delta t + e^{-\beta\Delta t}\,\E[Q(s_{t+1},a_{t+1})\mid
s_t,a_t]$, has the same structure.)

\subsection{Taylor-expanding the expectation}
Let $\Delta s := s_{t+1} - s_t$. Expanding $V(s_{t+1})$ about $s_t$ to
second order,
\begin{equation}
\label{eq:taylor}
V(s_{t+1}) \;\approx\; V(s_t)
\;+\; \Delta s^\top \grads V(s_t)
\;+\; \tfrac12\, \Delta s^\top \hess V(s_t)\, \Delta s .
\end{equation}
Taking the conditional expectation and introducing the (data-estimable)
drift and diffusion statistics
\begin{equation}
\label{eq:mu-sigma}
\bm{\mu} \;:=\; \frac{\E[\Delta s \mid s_t]}{\Delta t},
\qquad
\bm{\Sigma} \;:=\; \frac{\E[\Delta s\, \Delta s^\top \mid s_t]}{\Delta t},
\end{equation}
substituting \eqref{eq:taylor} into \eqref{eq:ct-value}, expanding
$e^{-\beta\Delta t} \approx 1 - \beta\Delta t$, and collecting the
first-order terms in $\Delta t$, we obtain the core PhiBE relation
\begin{equation}
\label{eq:phibe}
\boxed{\;
\beta\, V(s) \;=\; r
\;+\; \bm{\mu}^\top \grads V(s)
\;+\; \tfrac12\,\tr\!\big(\bm{\Sigma}\, \hess V(s)\big).
\;}
\end{equation}
Equation \eqref{eq:phibe} is the first point at which the three objects
$V$, $\grads V$, and $\hess V$ appear \emph{together}.

\section{The learning objective}

Define the per-state PhiBE residual of a parametric value function
$V_\theta$:
\begin{equation}
\label{eq:residual}
\Res_\theta(s) \;:=\;
\beta\, V_\theta(s) \;-\; r
\;-\; \bm{\mu}^\top \grads V_\theta(s)
\;-\; \tfrac12\,\tr\!\big(\bm{\Sigma}\, \hess V_\theta(s)\big).
\end{equation}
The objective to be minimized is the mean-squared residual over the data
distribution $\D$:
\begin{equation}
\label{eq:objective}
\boxed{\;
f(\theta) \;=\; \E_{s\sim\D}\!\Big[\, \Res_\theta(s)^2 \,\Big]
\;=\; \E_{s}\!\Big[\big(\beta V_\theta - r
- \bm{\mu}^\top \grads V_\theta
- \tfrac12 \tr(\bm{\Sigma}\, \hess V_\theta)\big)^2\Big].
\;}
\end{equation}

\paragraph{The three terms.}
Writing $d$ for the state dimension:
\begin{center}
\begin{tabular}{lll}
\hline
Term & Object & Shape \\
\hline
First  & $V_\theta(s)$        & scalar (coefficient $\beta$) \\
Second & $\grads V_\theta(s)$ & vector in $\R^{d}$ (contracted with $\bm{\mu}$) \\
Third  & $\hess V_\theta(s)$  & matrix in $\R^{d\times d}$ (trace against $\bm{\Sigma}$) \\
\hline
\end{tabular}
\end{center}
Crucially, $\bm{\mu}$ and $\bm{\Sigma}$ are estimated from data and carry
\emph{no} dependence on $\theta$; the only $\theta$-dependent objects are
$V_\theta$, $\grads V_\theta$, and $\hess V_\theta$.

\paragraph{Contrast with standard RL.}
A standard RL objective has the form
$f\big(V_\theta(s)\big) = \E[(V_\theta(s) - y)^2]$, where the target $y$ is
treated as a constant; it depends on $\theta$ only through the network's
\emph{output value}. Our objective \eqref{eq:objective} has the form
\begin{equation}
\label{eq:contrast}
\underbrace{f\big(V_\theta\big)}_{\text{standard RL}}
\qquad\longrightarrow\qquad
\underbrace{f\big(V_\theta,\ \grads V_\theta,\ \hess V_\theta\big)}_{\text{PhiBE}} ,
\end{equation}
i.e.\ the loss depends on $\theta$ through the output value
\emph{and} its first and second derivatives with respect to the input.

\section{Why a standard RL optimizer cannot minimize \eqref{eq:objective}}

Standard RL optimization (SGD/Adam with semi-gradient bootstrapping) is
built on the assumption that the loss depends on $\theta$ \emph{only}
through the output value. Objective \eqref{eq:objective} violates this
assumption, and three distinct difficulties follow.

\subsection*{Difficulty 1: mixed parameter--input higher-order derivatives}
Minimizing $f$ requires $\grad_\theta f$. By the chain rule,
\begin{equation}
\grad_\theta f \;=\; 2\,\E\big[\Res_\theta(s)\,\grad_\theta \Res_\theta(s)\big],
\end{equation}
and $\grad_\theta \Res_\theta$ contains
\begin{equation}
\grad_\theta\!\big(\grads V_\theta\big)
\qquad\text{and}\qquad
\grad_\theta\!\big(\hess V_\theta\big).
\end{equation}
A standard critic update needs only $\grad_\theta V_\theta$---a single
backward pass. Here we additionally need $\grad_\theta\grads V_\theta$ (a
mixed second derivative: differentiate once w.r.t.\ the input, once w.r.t.\
the parameters) and $\grad_\theta\hess V_\theta$ (a mixed third derivative).
These objects do not appear anywhere in a standard RL pipeline; computing
them requires nested differentiation through the network at every step,
which is expensive and numerically delicate.

\subsection*{Difficulty 2: the derivative terms are uncontrolled}
A standard objective constrains the network's \emph{output value} directly.
But nothing in \eqref{eq:objective} directly constrains the magnitude of
$\grads V_\theta$ (the input slope) or $\hess V_\theta$ (the input
curvature). A neural network can produce reasonable values while having
arbitrarily steep slopes. Consequently the term
$\bm{\mu}^\top \grads V_\theta$ entering the bootstrapped target can grow
without bound during training, driving the optimization to diverge.
A standard optimizer presumes the loss terms are well behaved; that
presumption fails here.

\subsection*{Difficulty 3: the second-order term is ill-posed for common networks}
The third term depends entirely on the second-order differentiability of
the network:
\begin{itemize}[leftmargin=1.5em]
\item For ReLU networks, $\hess V_\theta$ is (almost everywhere) zero, so
the third term vanishes identically and the expansion collapses to first
order, discarding the diffusion information the second-order term is meant
to carry.
\item For smooth activations (tanh/GELU), $\hess V_\theta$ exists but is
typically oscillatory and difficult to control in magnitude, making the
optimization less stable.
\end{itemize}

\subsection*{Summary of the difficulty}
The root cause is structural: a standard optimizer assumes the loss
depends on $\theta$ \emph{only through the output value}, whereas
\eqref{eq:objective} depends on $\theta$ through the \emph{value, the input
gradient, and the input Hessian}. This forces mixed higher-order
derivatives (Difficulty~1), exposes unconstrained quantities to the
bootstrapping dynamics (Difficulty~2), and renders the second-order term
ill-posed under common network choices (Difficulty~3).

\section{Next step}
The objective \eqref{eq:objective} is now stated precisely, and the reasons
a standard RL optimizer fails to minimize it are understood. The next stage
is to identify a suitable optimization method that can minimize an objective
depending jointly on $V_\theta$, $\grads V_\theta$, and $\hess V_\theta$.
We would value the advisor's guidance on what form this optimizer should
take before committing to an implementation.

\end{document}
